\section{A Field Built From a Title You Do Not Store}
\label{sec:ch09-a-field-built-from-a-title}

\index{requires@\code{"@requires}!the problem it solves|(}%
\code{averageScore} and \code{ratingCount} need nothing but the key. A third
field on the Ratings service needs something the service does not have. The
url where an attendee leaves a score is built from the session's title, and the
title is the Sessions service's:

\begin{minted}{csharp}
    [Requires("title")]
    public static string GetFeedbackUrl([Parent] Session session)
        => "https://feedback.example/2026/" + Slug(session.Title);
\end{minted}

\index{external@\code{"@external}!declares a field you do not resolve}%
For that to be legal the type has to declare \code{title}, and declare that it
is somebody else's:

\begin{minted}{csharp}
    [External]
    public static string GetTitle([Parent] Session session)
        => session.Title;
\end{minted}

\index{external@\code{"@external}!the resolver nobody calls}%
That method has a body and nothing ever calls it. \code{@external} says this
field is resolved elsewhere, so the router never routes a request for
\code{Session.title} here; the resolver exists because C\# wants a member
behind the field, and it returns whatever the reference resolver put in the
record. Writing a real body for a field marked external is a small dishonesty
that the language forces and the schema corrects.

\index{directives!where each attribute is allowed to go}%
Where the two attributes go is not a preference. \code{[Requires]} and
\code{[Provides]} derive from a base whose declared usage is methods and
properties only, so putting either on the type class does not compile:

\begin{minted}[fontsize=\footnotesize,breakanywhere]{text}
error CS0592: Attribute 'Requires' is not valid on this declaration type. It is only valid on 'method, property, indexer' declarations.
\end{minted}

\index{Key@\code{[Key]}!goes on the class, unlike these two}%
Which is the opposite of \code{[Key]} and \code{[ReferenceResolver]}, both of
which go on the class and neither of which is legal on a method.
\code{[External]} is allowed in both places. There is a rule underneath the
inconsistency and it is the same one the directives themselves follow: a
directive that describes a type goes on the type, and a directive that
describes one field goes on the field.

\subsection{What the Specification Says It Does}

\index{requires@\code{"@requires}!the normative sentence}%
Apollo's own reference is unusually direct about the mechanism, and the clause
worth reading twice is the last one. \code{@requires}
\enquote{tells the router that it needs to fetch the values of those externally
defined fields first, even if the original client query didn't
request them}~\autocite{apollodirectives}. It also states the precondition this
section has already met: a subgraph that requires a field must declare that
field and mark it external, or composition fails.

\index{requires@\code{"@requires}!fetches a field nobody asked for}%
That second clause is testable, so test it. Ask for the url and not the title:

\begin{minted}{graphql}
{ sessions(first: 10) { nodes { feedbackUrl } } }
\end{minted}

\begin{minted}[fontsize=\footnotesize,breakanywhere]{json}
{"data":{"sessions":{"nodes":[
  {"feedbackUrl":"https://feedback.example/2026/schemas-that-outlive-their-authors"},
  {"feedbackUrl":"https://feedback.example/2026/reading-a-query-plan"},
  {"feedbackUrl":"https://feedback.example/2026/paging-without-offsets"},
  {"feedbackUrl":"https://feedback.example/2026/the-bank-that-split-its-graph"}]}}}
\end{minted}

\index{requires@\code{"@requires}!the client pays for a field it never sees}%
No title anywhere in that response, and the Sessions service was asked for one
anyway: the request costs it a statement, and the plan's first fetch reads
\code{\{ sessions \{ nodes \{ title \_\_typename id \} \} \}}. A client asking
for one field caused a second field to be fetched, transmitted between two
services, and thrown away. That is the trade \code{@requires} makes, and it is
worth knowing before you put one on a field whose requirement is expensive.

\index{requires@\code{"@requires}!costs the requiring service nothing here}%
The Ratings service, meanwhile, issues no statement at all. Everything
\code{feedbackUrl} needs arrived in the representation.

\subsection{The Representation Grows a Second Entry}

\index{query plan!representations carry more than the key}%
Section~\ref{sec:ch09-the-first-query-across-two} read a
\code{representations} array with one entry in it, the key. With a
\code{@requires} in play there are two:

\begin{minted}[fontsize=\footnotesize]{json}
"representations": [
  { "kind": "@requires",
    "typeName": "Session",
    "fieldName": "feedbackUrl",
    "fragment": "fragment Requires_for_feedbackUrl on Session {\n    title\n}" },
  { "kind": "@key",
    "typeName": "Session",
    "fragment": "fragment Key on Session {\n    __typename\n    id\n}" }
]
\end{minted}

\index{requires@\code{"@requires}!is attached to a field, not a type}%
The first entry names the field that caused it, which is the difference between
this directive and the other two in the chapter. A key is a property of a type
and \code{@external} marks a field as somebody else's; a requirement belongs to
one field, and the plan carries it under that field's name.

\index{router.json@\code{router.json}!the compiled form of all three directives}%
All three of this chapter's directives compile into the execution config, each
into something small enough to read. The Ratings datasource is the only one
carrying a \code{requires} array:

\begin{minted}{json}
{ "typeName": "Session",
  "fieldName": "feedbackUrl",
  "selectionSet": "title" }
\end{minted}

\index{external@\code{"@external}!is a separate list in the config}%
And \code{@external} is a list beside the fields the subgraph actually
resolves, on the same datasource:

\begin{minted}[fontsize=\footnotesize]{json}
{ "typeName": "Session",
  "fieldNames": ["averageScore", "ratingCount", "feedbackUrl", "id"],
  "externalFieldNames": ["title"] }
\end{minted}

\index{composition!the directives are gone from what a client reads}%
Both of those are in the part of the config the router uses to plan. Neither
appears in the schema a client is served, along with \code{@key},
\code{@shareable} and the word \code{resolvable}, which is
chapter~\ref{ch:composition}'s point about composition holding for three new
directives: they do their work while the documents are merged and are gone from
what anyone downstream sees.

\subsection{A Field Set Nothing Checks}

\index{requires@\code{"@requires}!its argument is barely validated}%
One thing to know before you trust that string. The field set inside
\code{[Requires]} is parsed at schema build, so a syntax error stops the
service:

\begin{minted}[fontsize=\footnotesize,breakanywhere]{text}
Expected a `Name`-token, but found a `Bang`-token.
  at HotChocolate.ApolloFederation.Types.FieldSetType.ParseSelectionSet(String s)
  at HotChocolate.ApolloFederation.Types.RequiresDirective..ctor(String fields)
\end{minted}

\index{requires@\code{"@requires}!a typo reaches the exported schema}%
But a field set that parses and names a field that does not exist builds
cleanly and exports cleanly. \code{[Requires("doesNotExist")]} on a type with no
such field reaches the schema document as
\code{@requires(fields: "doesNotExist")}, untouched. The string is checked for
being a selection set and not for selecting anything, so the first thing that
notices a misspelled field name is downstream of the service that shipped it.
\index{requires@\code{"@requires}!the problem it solves|)}%
